\chapter*{Revisão Bibliográfica}
\addcontentsline{toc}{chapter}{Revisão Bibliográfica}

\section*{Visão Computacional Aplicada e Sensoriamento}
\addcontentsline{toc}{section}{Visão Computacional Aplicada e Sensoriamento}

\citeonline{BradskiKaehler2008} Definem visão computacional como a transformação de dados provenientes de uma câmera estática ou de vídeo em uma decisão ou nova representação.

\citeonline{Szeliski2021} Complementa esse conceito dizendo que a visão computacional pode ser definida como a área responsável por inferir e reconstruir propriedades do mundo físico a partir de imagens ou sequências de imagens, buscando estimar atributos como forma, iluminação, movimento e distribuição de cores. Embora se trate de um problema computacionalmente complexo, essa área já apresenta ampla maturidade tecnológica, sendo aplicada em diversos contextos do mundo real, como reconhecimento óptico de caracteres, inspeção industrial automatizada, logística e manipulação robótica em armazéns, diagnóstico por imagens médicas, reconstrução tridimensional por fotogrametria, efeitos visuais na indústria cinematográfica, vigilância inteligente, captura de movimento e sistemas biométricos. A diversidade dessas aplicações evidencia o papel central da visão computacional como elemento fundamental para sistemas autônomos e inteligentes.

As aplicações práticas de visão computacional demonstram versatilidade em diferentes setores. No monitoramento biológico, \citeonline{Cardoso2020} validou métodos eficazes para a contagem automatizada de espécimes em ambientes agropecuários, embora destaque a sensibilidade à iluminação como um desafio persistente. No campo da interação e acessibilidade, \citeonline{Cunha2022} comprovou a viabilidade de sistemas de rastreamento ocular de baixo custo utilizando hardware convencional, reforçando a tendência de democratização de soluções de VC que antes exigiam sensores de alto custo.

No âmbito destas aplicações, permeiam a literatura atualmente interseções da visão computacional com o xadrez. O trabalho de \citeonline{Lemevs2025Computer} apresenta como essa tecnologia pode ser eficaz na automação do reconhecimento de tabuleiros e peças para auxiliar a análise matemática de partidas ou permitir jogos contra computadores sob mesas físicas de xadrez. Provas de conceito dessa proposta são observáveis por exemplo na implementação de sistemas de rastreio de jogos de xadrez em MATLAB por \citeonline{Koray2016Computer} e em Python por \citeonline{Bugarin2024Real}, que se vale de ferramentas software livre e de código aberto como motor lógico do sistema.

\section*{Robótica e Controle Assistido por Visão}
\addcontentsline{toc}{section}{Robótica e Controle Assistido por Visão}

A integração entre percepção visual e controle robótico é explorada tanto em sistemas aéreos quanto terrestres e industriais. \citeonline{Menezes2013} contribui com a análise de controladores alternativos para servovisão em robôs aéreos, apontando melhorias de desempenho em cenários de voo. No âmbito da robótica educativa e de pesquisa, \citeonline{Rezeck2019} estabeleceu bases para plataformas abertas de baixo custo, enquanto \citeonline{Bertoncini2022} e \citeonline{Santos2022} avançaram na automação aplicada: o primeiro focando em centrais de controle para navegação autônoma e o segundo na precisão de manipuladores robóticos industriais para tarefas de \textit{pick-and-place}.

A automação integrada ao jogo de xadrez é também emergente neste contexto. Os autores \citeonline{Sarker2015Wizard} propuseram o \textit{Wizard chess}: Um sistema autônomo que utiliza um braço robótico articulado aliado à um tabuleiro magnético e algoritmos de visão computacional para permitir liberdade de ação completa do sistema sob o jogo. O \textit{Wizard chess} se destaca por seu baixo custo no cenário de uso pessoal ou educacional além de sua arquitetura pouco complexa, demonstrando a eficiência da tecnologia robótica em seu estado-da-arte.

\section*{Infraestrutura IoT, Segurança e Inteligência de Dados}
\addcontentsline{toc}{section}{Infraestrutura IoT, Segurança e Inteligência de Dados}

A sustentação dessas tecnologias depende de redes eficientes e seguras. \citeonline{Pereira2022} demonstra como o Aprendizado Profundo pode ser aplicado na classificação de tráfego em redes definidas por software (SDN) integradas à IoT, essencial para o desenvolvimento da Indústria 4.0 preconizada por \citeonline{Schwab2016}. Paralelamente, a eficiência energética e a segurança surgem como pilares críticos; enquanto \citeonline{GuimaraesDaSilva2023} oferece um comparativo rigoroso de modelos de DL para previsão de consumo de energia em dispositivos IoT, \citeonline{BritoGuimaraes2023} aborda a vulnerabilidade desses sistemas, utilizando hardware de borda (\textit{Edge Computing}) para a detecção de \textit{botnets} em tempo real.

\section*{Síntese da Revisão Bibliográfica}
\addcontentsline{toc}{section}{Síntese da Revisão Bibliográfica}

De forma convergente, a literatura indica que, embora a robótica moderna dependa fortemente de sistemas de percepção visual e de modelos avançados de aprendizado profundo, sua aplicação prática em plataformas embarcadas e de baixo custo ainda enfrenta limitações significativas relacionadas ao poder computacional, consumo energético e latência. Conforme apontado por \citeonline{Szeliski2021} e \citeonline{BradskiKaehler2008}, a complexidade inerente aos algoritmos de visão computacional cresce substancialmente à medida que se busca maior robustez e autonomia. Paralelamente, estudos recentes em IoT e \textit{Edge Computing} \citeonline{Pereira2022,BritoGuimaraes2023} reforçam a necessidade de arquiteturas distribuídas e otimizadas, capazes de executar inferência local sem comprometer a responsividade do sistema. Nesse contexto, evidencia-se uma lacuna de pesquisa relacionada à integração eficiente entre visão computacional, inteligência artificial embarcada e controle robótico em manipuladores antropomórficos de baixo custo, especialmente quando se busca autonomia operacional em tempo real — lacuna esta que o presente projeto se propõe a investigar.